{\fontsize{10}{15}\selectfont
译文采用无韵自由诗，按所据拉丁正文逐行对应排列；语意有时跨越相邻诗行承接。正文左侧页边每五行标示原诗行号，注释置于页脚。本译本以 Gian Biagio Conte 编订的 \href{https://doi.org/10.1515/9783110649031}{\textit{Aeneis: Editio altera}}（De Gruyter，2019）第二版为主要拉丁底本，正文及行序原则上从其文本；疑难处参考 Servius 古注、现代各卷注释及相关研究。个别存在异文或解释分歧之处，酌据语法、上下文及注释传统裁决，并于确有必要时在脚注中说明。书末“专名与术语表”收录阅读中反复出现、意义容易分散或较易混淆的核心术语、称号、人物、神祇与地名。
\par
注释中“Conte 2019”指上述拉丁校勘本；未另署出处的“校勘记”指该版拉丁正文页下记录传本异文、学者校订及删疑意见的部分。
\par
来源标记说明：古代作者在注释中使用中文名；“塞尔维乌斯”（Servius）指\href{https://la.wikisource.org/wiki/In_Vergilii_Aeneidem_commentarii}{《埃涅阿斯纪评注》}（Thilo 本，Wikisource 数字文本），括号内为所评卷、行，依本稿诗行编号。马克罗比乌斯括号内的三组数字指\href{https://penelope.uchicago.edu/Thayer/L/Roman/Texts/Macrobius/Saturnalia/6*.html}{《农神节》}的卷、章、节，在线文本见 LacusCurtius。现代文献首次引用时注明作品，后续使用作者与年份，并标明所评诗行或讨论位置；转引注明所据来源。经 Conte 校勘记转述的学者意见保留这一转述关系。
\par}
